% GATE CS2 2026: formulae used in the solutions
% Included from main.tex with % GATE CS2 2026: formulae used in the solutions
% Included from main.tex with % GATE CS2 2026: formulae used in the solutions
% Included from main.tex with % GATE CS2 2026: formulae used in the solutions
% Included from main.tex with \input{chapters/cs2-2026_formulae.tex}

% Unique link targets for equations, figures and tables. Numbers restart in
% every item (\numberwithin{...}{enumi} in main.tex), so without the
% subsection in the name, e.g. equation 1 of item 1 in Formulae and in
% Problems would get the same hyperref target and links could jump to the
% wrong place ("destination with the same identifier" warnings).
\renewcommand{\theHequation}{\thesubsection.\arabic{enumi}.\arabic{equation}}
\renewcommand{\theHfigure}{\thesubsection.\arabic{enumi}.\arabic{figure}}
\renewcommand{\theHtable}{\thesubsection.\arabic{enumi}.\arabic{table}}

\begin{enumerate}[label=\thesubsection.\arabic*, ref=\thesubsection.\theenumi]

%---------------------------------------------------------------------
\item The Z-transform of the PMF of an integer-valued random variable $X$
is defined as
\begin{align}
M_X(z) = \sum_{k=-\infty}^{\infty} p_X(k)\, z^{-k} = \mean{z^{-X}}
\label{eq:ztrans-def}
\end{align}

%---------------------------------------------------------------------
\item If $X_1, \dots, X_n$ are independent and $Y = X_1 + \dots + X_n$, then
\begin{align}
M_Y(z) = \prod_{i=1}^{n} M_{X_i}(z)
\label{eq:ztrans-prod}
\end{align}

%---------------------------------------------------------------------
\item For a Bernoulli random variable $X$ with $\pr{X = 1} = p$ and
$\pr{X = 0} = q = 1 - p$,
\begin{align}
M_X(z) = q + p\, z^{-1}
\label{eq:ztrans-bern}
\end{align}

%---------------------------------------------------------------------
\item If $Y$ is the sum of $n$ independent Bernoulli random variables, each
with parameter $p$, then
\begin{align}
p_Y(k) = \binom{n}{k} p^k q^{n-k}, \quad k = 0, 1, \dots, n
\label{eq:binom-pmf}
\end{align}

%---------------------------------------------------------------------
\item The joint Z-transform of two integer-valued random variables $X$ and
$Y$ is
\begin{align}
M_{X,Y}\brak{z_1, z_2} = \mean{z_1^{-X} z_2^{-Y}}
= \sum_{j}\sum_{k} \pr{X = j, Y = k}\, z_1^{-j} z_2^{-k}
\label{eq:ztrans-joint}
\end{align}
so $\pr{X = j, Y = k}$ is the coefficient of $z_1^{-j} z_2^{-k}$. If
$\brak{X, Y} = \sum_i \brak{X_i, Y_i}$ with the pairs $\brak{X_i, Y_i}$
independent, then
\begin{align}
M_{X,Y}\brak{z_1, z_2} = \prod_{i} M_{X_i,Y_i}\brak{z_1, z_2}
\label{eq:ztrans-joint-prod}
\end{align}

%---------------------------------------------------------------------
\item For a fair die $X$, uniform on $\cbrak{1, \dots, 6}$,
\begin{align}
M_X(z) = \frac{1}{6}\sum_{k=1}^{6} z^{-k}
\label{eq:ztrans-die}
\end{align}
If $S = X_1 + X_2$ for two independent fair dice,
\begin{align}
p_S(k) = \frac{6 - \abs{k - 7}}{36}, \quad k = 2, \dots, 12
\label{eq:dice-sum-pmf}
\end{align}

%---------------------------------------------------------------------
\item A random variable $X$ is Gaussian (normal), $X \sim \gauss{\mu}{\sigma^2}$,
if its PDF is
\begin{align}
f(x) = \frac{1}{\sigma\sqrt{2\pi}}
\exp\brak{-\frac{\brak{x - \mu}^2}{2\sigma^2}}, \quad
x \in \brak{-\infty, \infty}
\label{eq:gauss-pdf}
\end{align}

%---------------------------------------------------------------------
\item Other standard distributions:
\begin{enumerate}[label=\alph*)]
\item An exponential random variable with rate $\lambda > 0$ is continuous,
takes only non-negative values, and has PDF
\begin{align}
f(x) =
\begin{cases}
\lambda e^{-\lambda x} & x \ge 0 \\
0 & x < 0
\end{cases}
\label{eq:exp-pdf}
\end{align}
\item A Poisson random variable $X$ with parameter $\lambda > 0$ is
discrete. It counts the number of events in a fixed interval when events
occur independently at an average rate $\lambda$, so it takes only the
values $0, 1, 2, \dots$, with PMF
\begin{align}
\pr{X = k} = \frac{\lambda^{k} e^{-\lambda}}{k!}, \quad k = 0, 1, 2, \dots
\label{eq:poisson-pmf}
\end{align}
Being discrete, it has a PMF, not a PDF over the real line.
\item A uniform random variable on $[a, b]$ is continuous with PDF
\begin{align}
f(x) =
\begin{cases}
\frac{1}{b - a} & a \le x \le b \\
0 & \text{otherwise}
\end{cases}
\label{eq:unif-pdf}
\end{align}
which is constant on $[a, b]$ and zero outside it.
\end{enumerate}

%---------------------------------------------------------------------
\item Conditional probability:
\begin{align}
\pr{A \mid B} = \frac{\pr{AB}}{\pr{B}}
\label{eq:cond-prob}
\end{align}

%---------------------------------------------------------------------
\item If $g$ is an odd function, $g(-x) = -g(x)$, then
\begin{align}
\int_{-a}^{a} g(x)\, dx = 0
\label{eq:odd-int}
\end{align}

%---------------------------------------------------------------------
\item For $m \neq -1$,
\begin{align}
\int x^{m}\, dx = \frac{x^{m+1}}{m+1} + C
\label{eq:power-rule}
\end{align}

%---------------------------------------------------------------------
\item The sum of the first $n$ natural numbers is
\begin{align}
\sum_{k=1}^{n} k = \frac{n(n+1)}{2}
\label{eq:arith-sum}
\end{align}

%---------------------------------------------------------------------
\item Asymptotic notation. For positive functions $f(n)$ and $g(n)$,
\begin{align}
f(n) = o\brak{g(n)} &\iff \lim_{n \to \infty} \frac{f(n)}{g(n)} = 0
\label{eq:little-o} \\
f(n) = O\brak{g(n)} &\iff f(n) \le c\, g(n) \text{ for large } n
\label{eq:big-o} \\
f(n) = \Theta\brak{g(n)} &\iff c_1 g(n) \le f(n) \le c_2 g(n)
\text{ for large } n
\label{eq:theta}
\end{align}
for some constants $c, c_1, c_2 > 0$. In words, $o$ (little-o) means $f$ grows \emph{strictly slower} than $g$:
however large a multiple of $f$ is taken, $g$ eventually overtakes it.
$O$ (big-O) means $f$ grows \emph{no faster} than $g$, possibly at the
same rate. $\Theta$ means $f$ and $g$ grow at the \emph{same} rate, up to
constant factors. So $o$ is the strict version of $O$, just as $<$ is the
strict version of $\le$.

%---------------------------------------------------------------------
\item For any constants $\epsilon > 0$ and $0 \le a < b$,
\begin{align}
\log(n) = o\brak{n^{\epsilon}}, \qquad n^{a} = o\brak{n^{b}}
\label{eq:growth}
\end{align}

%---------------------------------------------------------------------
\item
\begin{align}
\log(n!) = \Theta\brak{n \log(n)}
\label{eq:log-fact}
\end{align}

%---------------------------------------------------------------------
\item Master theorem. For $T(n) = a\,T\brak{\frac{n}{b}} + f(n)$ with
$a \ge 1$, $b > 1$,
\begin{align}
T(n) =
\begin{cases}
\Theta\brak{n^{\log_b a}} & f(n) = O\brak{n^{\log_b a - \epsilon}}
\text{ for some } \epsilon > 0 \\
\Theta\brak{n^{\log_b a}\log(n)} & f(n) = \Theta\brak{n^{\log_b a}}
\end{cases}
\label{eq:master}
\end{align}

%---------------------------------------------------------------------
\item The (one-sided) Z-transform of a sequence $x[k]$, $k \ge 0$, is
\begin{align}
X(z) = \sum_{k=0}^{\infty} x[k]\, z^{-k}
\label{eq:zseq-def}
\end{align}
If $x[-1] = 0$, a delay by one step multiplies the transform by $z^{-1}$:
\begin{align}
x[k-1] \system{Z} z^{-1} X(z)
\label{eq:zseq-shift}
\end{align}

%---------------------------------------------------------------------
\item For a constant $a$ and $\abs{z} > \abs{a}$,
\begin{align}
a^{k} &\system{Z} \frac{1}{1 - a z^{-1}}
\label{eq:zseq-geo} \\
(k+1)\, a^{k} &\system{Z} \frac{1}{\brak{1 - a z^{-1}}^{2}}
\label{eq:zseq-geo2} \\
\frac{(k+1)(k+2)}{2} &\system{Z} \frac{1}{\brak{1 - z^{-1}}^{3}}
\label{eq:zseq-geo3}
\end{align}

%---------------------------------------------------------------------
\item A relation $R$ on a set $A$ is
\begin{align}
\text{reflexive if } & (x, x) \in R \ \ \forall x \in A
\label{eq:rel-ref} \\
\text{symmetric if } & (x, y) \in R \implies (y, x) \in R
\label{eq:rel-sym} \\
\text{transitive if } & (x, y), (y, z) \in R \implies (x, z) \in R
\label{eq:rel-trans} \\
\text{antisymmetric if } & (x, y), (y, x) \in R \implies x = y
\label{eq:rel-anti}
\end{align}

%---------------------------------------------------------------------
\item The characteristic equation for a matrix $\vec{A}$ is
\begin{align}
f\brak{\lambda} = \mydet{\vec{A} - \lambda\vec{I}} = 0
\label{eq:chareq}
\end{align}

%---------------------------------------------------------------------
\item Cayley-Hamilton theorem:
\begin{align}
f\brak{\lambda} = f\brak{\vec{A}} = \vec{0}
\label{eq:cayley}
\end{align}

%---------------------------------------------------------------------
\item Solving $\vec{A}\vec{x} = \vec{c}$ by row reduction of the augmented
matrix $\myvec{\vec{A} & \vec{c}}$. Elementary row operations, such as
\begin{align}
R_2 \leftarrow \alpha R_2 + \beta R_1, \quad \alpha \ne 0
\label{eq:row-op}
\end{align}
do not change the solutions. After reduction:
\begin{enumerate}[label=\alph*)]
\item a row $\myvec{0 & \dots & 0 & | & 0}$ leaves a variable free, giving
infinitely many solutions;
\item a row $\myvec{0 & \dots & 0 & | & d}$ with $d \ne 0$ means there is
no solution.
\end{enumerate}
\label{fm:row-reduce}

%---------------------------------------------------------------------
\item For an $n \times n$ matrix $\vec{A}$ and a scalar $c$,
\begin{align}
\det\brak{c\vec{A}} = c^{n}\det\brak{\vec{A}}
\label{eq:det-scale}
\end{align}

%---------------------------------------------------------------------
\item In $n$-bit two's complement, the bits $b_{n-1} \dots b_1 b_0$
represent
\begin{align}
x = -b_{n-1} 2^{n-1} + \sum_{i=0}^{n-2} b_i 2^{i}
\label{eq:twos-val}
\end{align}
and the representable range is
\begin{align}
-2^{n-1} \le x \le 2^{n-1} - 1
\label{eq:twos-range}
\end{align}
An arithmetic operation overflows if its true result lies outside
\eqref{eq:twos-range}.

%---------------------------------------------------------------------
\item IEEE 754 single precision. A 32-bit word with sign bit $s$, 8-bit
exponent $E$ and 23-bit fraction $f$ represents (for $0 < E < 255$)
\begin{align}
x = \brak{-1}^{s} \times \brak{1.f}_2 \times 2^{E - 127}
\label{eq:ieee}
\end{align}

%---------------------------------------------------------------------
\item Cyclic Redundancy Check. For data bits $D(x)$ and a generator $G(x)$
of degree $r$, the check bits appended to the data are
\begin{align}
R(x) = x^{r} D(x) \bmod G(x)
\label{eq:crc}
\end{align}
computed with modulo-2 (XOR) arithmetic. In bits: append $r$ zeros to
the data and divide by the generator pattern.

%---------------------------------------------------------------------
\item Boolean algebra:
\begin{align}
X + X' &= 1 \label{eq:bool-comp} \\
XY + X'Y &= Y \label{eq:bool-comb}
\end{align}

%---------------------------------------------------------------------
\item An active-high $2:4$ decoder with inputs $A$ (MSB) and $B$ gives
\begin{align}
D_0 = A'B', \quad D_1 = A'B, \quad D_2 = AB', \quad D_3 = AB
\label{eq:decoder}
\end{align}
A $4$-to-$1$ multiplexer with select lines $S_1 S_0$ and inputs
$I_0, \dots, I_3$ gives
\begin{align}
Y = S_1'S_0' I_0 + S_1'S_0 I_1 + S_1 S_0' I_2 + S_1 S_0 I_3
\label{eq:mux}
\end{align}

%---------------------------------------------------------------------
\item The Jaccard coefficient of two sets $A$, $B$ is
\begin{align}
J(A, B) = \frac{\abs{A \cap B}}{\abs{A \cup B}}
\label{eq:jaccard}
\end{align}
A spanning tree of a connected graph with $n$ vertices has $n - 1$ edges.
The complete graph $K_n$ has $\frac{n(n-1)}{2}$ edges.

%---------------------------------------------------------------------
\item Hashing with chaining. A key $k$ is placed in the chain of slot
\begin{align}
h(k) = k \bmod m
\label{eq:hash}
\end{align}
for a table of $m$ slots. A chain's length is the number of keys that
hash to its slot.

%---------------------------------------------------------------------
\item Merging sorted arrays $L$ and $R$ outputs $L$ followed by $R$ if and
only if
\begin{align}
\max(L) \le \min(R)
\label{eq:merge-void}
\end{align}

%---------------------------------------------------------------------
\item A stack is last-in first-out: elements are popped in the reverse of
the order in which they were pushed. A queue is first-in first-out:
elements are dequeued in the order in which they were enqueued.
\label{fm:stack-queue}

%---------------------------------------------------------------------
\item Binary semaphore $S$:
\begin{align}
\text{Wait}(S) &: \text{block until } S = 1, \text{ then set } S = 0
\label{eq:sem-wait} \\
\text{Signal}(S) &: \text{set } S = 1 \text{ (waking one waiting process)}
\label{eq:sem-signal}
\end{align}
A binary semaphore initialised to $1$, with Wait before and Signal after a
block of code, lets only one process run that block at a time (mutual
exclusion).

%---------------------------------------------------------------------
\item C scope rules. A declaration in an inner block hides a declaration
of the same name in an outer scope. The same name cannot be declared twice
in the same scope. For any variable \texttt{x},
\begin{align}
\texttt{*(\&x)} = \texttt{x}
\label{eq:deref}
\end{align}
\label{fm:c-scope}

%---------------------------------------------------------------------
\item The second digit after the decimal point of $r \in (0, 1)$ is
\begin{align}
d_2(r) = \lfloor 100 r \rfloor \bmod 10
\label{eq:digit2}
\end{align}
so $d_2$ is constant on each interval
$\left[\frac{k}{100}, \frac{k+1}{100}\right)$.

\end{enumerate}


% Unique link targets for equations, figures and tables. Numbers restart in
% every item (\numberwithin{...}{enumi} in main.tex), so without the
% subsection in the name, e.g. equation 1 of item 1 in Formulae and in
% Problems would get the same hyperref target and links could jump to the
% wrong place ("destination with the same identifier" warnings).
\renewcommand{\theHequation}{\thesubsection.\arabic{enumi}.\arabic{equation}}
\renewcommand{\theHfigure}{\thesubsection.\arabic{enumi}.\arabic{figure}}
\renewcommand{\theHtable}{\thesubsection.\arabic{enumi}.\arabic{table}}

\begin{enumerate}[label=\thesubsection.\arabic*, ref=\thesubsection.\theenumi]

%---------------------------------------------------------------------
\item The Z-transform of the PMF of an integer-valued random variable $X$
is defined as
\begin{align}
M_X(z) = \sum_{k=-\infty}^{\infty} p_X(k)\, z^{-k} = \mean{z^{-X}}
\label{eq:ztrans-def}
\end{align}

%---------------------------------------------------------------------
\item If $X_1, \dots, X_n$ are independent and $Y = X_1 + \dots + X_n$, then
\begin{align}
M_Y(z) = \prod_{i=1}^{n} M_{X_i}(z)
\label{eq:ztrans-prod}
\end{align}

%---------------------------------------------------------------------
\item For a Bernoulli random variable $X$ with $\pr{X = 1} = p$ and
$\pr{X = 0} = q = 1 - p$,
\begin{align}
M_X(z) = q + p\, z^{-1}
\label{eq:ztrans-bern}
\end{align}

%---------------------------------------------------------------------
\item If $Y$ is the sum of $n$ independent Bernoulli random variables, each
with parameter $p$, then
\begin{align}
p_Y(k) = \binom{n}{k} p^k q^{n-k}, \quad k = 0, 1, \dots, n
\label{eq:binom-pmf}
\end{align}

%---------------------------------------------------------------------
\item The joint Z-transform of two integer-valued random variables $X$ and
$Y$ is
\begin{align}
M_{X,Y}\brak{z_1, z_2} = \mean{z_1^{-X} z_2^{-Y}}
= \sum_{j}\sum_{k} \pr{X = j, Y = k}\, z_1^{-j} z_2^{-k}
\label{eq:ztrans-joint}
\end{align}
so $\pr{X = j, Y = k}$ is the coefficient of $z_1^{-j} z_2^{-k}$. If
$\brak{X, Y} = \sum_i \brak{X_i, Y_i}$ with the pairs $\brak{X_i, Y_i}$
independent, then
\begin{align}
M_{X,Y}\brak{z_1, z_2} = \prod_{i} M_{X_i,Y_i}\brak{z_1, z_2}
\label{eq:ztrans-joint-prod}
\end{align}

%---------------------------------------------------------------------
\item For a fair die $X$, uniform on $\cbrak{1, \dots, 6}$,
\begin{align}
M_X(z) = \frac{1}{6}\sum_{k=1}^{6} z^{-k}
\label{eq:ztrans-die}
\end{align}
If $S = X_1 + X_2$ for two independent fair dice,
\begin{align}
p_S(k) = \frac{6 - \abs{k - 7}}{36}, \quad k = 2, \dots, 12
\label{eq:dice-sum-pmf}
\end{align}

%---------------------------------------------------------------------
\item A random variable $X$ is Gaussian (normal), $X \sim \gauss{\mu}{\sigma^2}$,
if its PDF is
\begin{align}
f(x) = \frac{1}{\sigma\sqrt{2\pi}}
\exp\brak{-\frac{\brak{x - \mu}^2}{2\sigma^2}}, \quad
x \in \brak{-\infty, \infty}
\label{eq:gauss-pdf}
\end{align}

%---------------------------------------------------------------------
\item Other standard distributions:
\begin{enumerate}[label=\alph*)]
\item An exponential random variable with rate $\lambda > 0$ is continuous,
takes only non-negative values, and has PDF
\begin{align}
f(x) =
\begin{cases}
\lambda e^{-\lambda x} & x \ge 0 \\
0 & x < 0
\end{cases}
\label{eq:exp-pdf}
\end{align}
\item A Poisson random variable $X$ with parameter $\lambda > 0$ is
discrete. It counts the number of events in a fixed interval when events
occur independently at an average rate $\lambda$, so it takes only the
values $0, 1, 2, \dots$, with PMF
\begin{align}
\pr{X = k} = \frac{\lambda^{k} e^{-\lambda}}{k!}, \quad k = 0, 1, 2, \dots
\label{eq:poisson-pmf}
\end{align}
Being discrete, it has a PMF, not a PDF over the real line.
\item A uniform random variable on $[a, b]$ is continuous with PDF
\begin{align}
f(x) =
\begin{cases}
\frac{1}{b - a} & a \le x \le b \\
0 & \text{otherwise}
\end{cases}
\label{eq:unif-pdf}
\end{align}
which is constant on $[a, b]$ and zero outside it.
\end{enumerate}

%---------------------------------------------------------------------
\item Conditional probability:
\begin{align}
\pr{A \mid B} = \frac{\pr{AB}}{\pr{B}}
\label{eq:cond-prob}
\end{align}

%---------------------------------------------------------------------
\item If $g$ is an odd function, $g(-x) = -g(x)$, then
\begin{align}
\int_{-a}^{a} g(x)\, dx = 0
\label{eq:odd-int}
\end{align}

%---------------------------------------------------------------------
\item For $m \neq -1$,
\begin{align}
\int x^{m}\, dx = \frac{x^{m+1}}{m+1} + C
\label{eq:power-rule}
\end{align}

%---------------------------------------------------------------------
\item The sum of the first $n$ natural numbers is
\begin{align}
\sum_{k=1}^{n} k = \frac{n(n+1)}{2}
\label{eq:arith-sum}
\end{align}

%---------------------------------------------------------------------
\item Asymptotic notation. For positive functions $f(n)$ and $g(n)$,
\begin{align}
f(n) = o\brak{g(n)} &\iff \lim_{n \to \infty} \frac{f(n)}{g(n)} = 0
\label{eq:little-o} \\
f(n) = O\brak{g(n)} &\iff f(n) \le c\, g(n) \text{ for large } n
\label{eq:big-o} \\
f(n) = \Theta\brak{g(n)} &\iff c_1 g(n) \le f(n) \le c_2 g(n)
\text{ for large } n
\label{eq:theta}
\end{align}
for some constants $c, c_1, c_2 > 0$. In words, $o$ (little-o) means $f$ grows \emph{strictly slower} than $g$:
however large a multiple of $f$ is taken, $g$ eventually overtakes it.
$O$ (big-O) means $f$ grows \emph{no faster} than $g$, possibly at the
same rate. $\Theta$ means $f$ and $g$ grow at the \emph{same} rate, up to
constant factors. So $o$ is the strict version of $O$, just as $<$ is the
strict version of $\le$.

%---------------------------------------------------------------------
\item For any constants $\epsilon > 0$ and $0 \le a < b$,
\begin{align}
\log(n) = o\brak{n^{\epsilon}}, \qquad n^{a} = o\brak{n^{b}}
\label{eq:growth}
\end{align}

%---------------------------------------------------------------------
\item
\begin{align}
\log(n!) = \Theta\brak{n \log(n)}
\label{eq:log-fact}
\end{align}

%---------------------------------------------------------------------
\item Master theorem. For $T(n) = a\,T\brak{\frac{n}{b}} + f(n)$ with
$a \ge 1$, $b > 1$,
\begin{align}
T(n) =
\begin{cases}
\Theta\brak{n^{\log_b a}} & f(n) = O\brak{n^{\log_b a - \epsilon}}
\text{ for some } \epsilon > 0 \\
\Theta\brak{n^{\log_b a}\log(n)} & f(n) = \Theta\brak{n^{\log_b a}}
\end{cases}
\label{eq:master}
\end{align}

%---------------------------------------------------------------------
\item The (one-sided) Z-transform of a sequence $x[k]$, $k \ge 0$, is
\begin{align}
X(z) = \sum_{k=0}^{\infty} x[k]\, z^{-k}
\label{eq:zseq-def}
\end{align}
If $x[-1] = 0$, a delay by one step multiplies the transform by $z^{-1}$:
\begin{align}
x[k-1] \system{Z} z^{-1} X(z)
\label{eq:zseq-shift}
\end{align}

%---------------------------------------------------------------------
\item For a constant $a$ and $\abs{z} > \abs{a}$,
\begin{align}
a^{k} &\system{Z} \frac{1}{1 - a z^{-1}}
\label{eq:zseq-geo} \\
(k+1)\, a^{k} &\system{Z} \frac{1}{\brak{1 - a z^{-1}}^{2}}
\label{eq:zseq-geo2} \\
\frac{(k+1)(k+2)}{2} &\system{Z} \frac{1}{\brak{1 - z^{-1}}^{3}}
\label{eq:zseq-geo3}
\end{align}

%---------------------------------------------------------------------
\item A relation $R$ on a set $A$ is
\begin{align}
\text{reflexive if } & (x, x) \in R \ \ \forall x \in A
\label{eq:rel-ref} \\
\text{symmetric if } & (x, y) \in R \implies (y, x) \in R
\label{eq:rel-sym} \\
\text{transitive if } & (x, y), (y, z) \in R \implies (x, z) \in R
\label{eq:rel-trans} \\
\text{antisymmetric if } & (x, y), (y, x) \in R \implies x = y
\label{eq:rel-anti}
\end{align}

%---------------------------------------------------------------------
\item The characteristic equation for a matrix $\vec{A}$ is
\begin{align}
f\brak{\lambda} = \mydet{\vec{A} - \lambda\vec{I}} = 0
\label{eq:chareq}
\end{align}

%---------------------------------------------------------------------
\item Cayley-Hamilton theorem:
\begin{align}
f\brak{\lambda} = f\brak{\vec{A}} = \vec{0}
\label{eq:cayley}
\end{align}

%---------------------------------------------------------------------
\item Solving $\vec{A}\vec{x} = \vec{c}$ by row reduction of the augmented
matrix $\myvec{\vec{A} & \vec{c}}$. Elementary row operations, such as
\begin{align}
R_2 \leftarrow \alpha R_2 + \beta R_1, \quad \alpha \ne 0
\label{eq:row-op}
\end{align}
do not change the solutions. After reduction:
\begin{enumerate}[label=\alph*)]
\item a row $\myvec{0 & \dots & 0 & | & 0}$ leaves a variable free, giving
infinitely many solutions;
\item a row $\myvec{0 & \dots & 0 & | & d}$ with $d \ne 0$ means there is
no solution.
\end{enumerate}
\label{fm:row-reduce}

%---------------------------------------------------------------------
\item For an $n \times n$ matrix $\vec{A}$ and a scalar $c$,
\begin{align}
\det\brak{c\vec{A}} = c^{n}\det\brak{\vec{A}}
\label{eq:det-scale}
\end{align}

%---------------------------------------------------------------------
\item In $n$-bit two's complement, the bits $b_{n-1} \dots b_1 b_0$
represent
\begin{align}
x = -b_{n-1} 2^{n-1} + \sum_{i=0}^{n-2} b_i 2^{i}
\label{eq:twos-val}
\end{align}
and the representable range is
\begin{align}
-2^{n-1} \le x \le 2^{n-1} - 1
\label{eq:twos-range}
\end{align}
An arithmetic operation overflows if its true result lies outside
\eqref{eq:twos-range}.

%---------------------------------------------------------------------
\item IEEE 754 single precision. A 32-bit word with sign bit $s$, 8-bit
exponent $E$ and 23-bit fraction $f$ represents (for $0 < E < 255$)
\begin{align}
x = \brak{-1}^{s} \times \brak{1.f}_2 \times 2^{E - 127}
\label{eq:ieee}
\end{align}

%---------------------------------------------------------------------
\item Cyclic Redundancy Check. For data bits $D(x)$ and a generator $G(x)$
of degree $r$, the check bits appended to the data are
\begin{align}
R(x) = x^{r} D(x) \bmod G(x)
\label{eq:crc}
\end{align}
computed with modulo-2 (XOR) arithmetic. In bits: append $r$ zeros to
the data and divide by the generator pattern.

%---------------------------------------------------------------------
\item Boolean algebra:
\begin{align}
X + X' &= 1 \label{eq:bool-comp} \\
XY + X'Y &= Y \label{eq:bool-comb}
\end{align}

%---------------------------------------------------------------------
\item An active-high $2:4$ decoder with inputs $A$ (MSB) and $B$ gives
\begin{align}
D_0 = A'B', \quad D_1 = A'B, \quad D_2 = AB', \quad D_3 = AB
\label{eq:decoder}
\end{align}
A $4$-to-$1$ multiplexer with select lines $S_1 S_0$ and inputs
$I_0, \dots, I_3$ gives
\begin{align}
Y = S_1'S_0' I_0 + S_1'S_0 I_1 + S_1 S_0' I_2 + S_1 S_0 I_3
\label{eq:mux}
\end{align}

%---------------------------------------------------------------------
\item The Jaccard coefficient of two sets $A$, $B$ is
\begin{align}
J(A, B) = \frac{\abs{A \cap B}}{\abs{A \cup B}}
\label{eq:jaccard}
\end{align}
A spanning tree of a connected graph with $n$ vertices has $n - 1$ edges.
The complete graph $K_n$ has $\frac{n(n-1)}{2}$ edges.

%---------------------------------------------------------------------
\item Hashing with chaining. A key $k$ is placed in the chain of slot
\begin{align}
h(k) = k \bmod m
\label{eq:hash}
\end{align}
for a table of $m$ slots. A chain's length is the number of keys that
hash to its slot.

%---------------------------------------------------------------------
\item Merging sorted arrays $L$ and $R$ outputs $L$ followed by $R$ if and
only if
\begin{align}
\max(L) \le \min(R)
\label{eq:merge-void}
\end{align}

%---------------------------------------------------------------------
\item A stack is last-in first-out: elements are popped in the reverse of
the order in which they were pushed. A queue is first-in first-out:
elements are dequeued in the order in which they were enqueued.
\label{fm:stack-queue}

%---------------------------------------------------------------------
\item Binary semaphore $S$:
\begin{align}
\text{Wait}(S) &: \text{block until } S = 1, \text{ then set } S = 0
\label{eq:sem-wait} \\
\text{Signal}(S) &: \text{set } S = 1 \text{ (waking one waiting process)}
\label{eq:sem-signal}
\end{align}
A binary semaphore initialised to $1$, with Wait before and Signal after a
block of code, lets only one process run that block at a time (mutual
exclusion).

%---------------------------------------------------------------------
\item C scope rules. A declaration in an inner block hides a declaration
of the same name in an outer scope. The same name cannot be declared twice
in the same scope. For any variable \texttt{x},
\begin{align}
\texttt{*(\&x)} = \texttt{x}
\label{eq:deref}
\end{align}
\label{fm:c-scope}

%---------------------------------------------------------------------
\item The second digit after the decimal point of $r \in (0, 1)$ is
\begin{align}
d_2(r) = \lfloor 100 r \rfloor \bmod 10
\label{eq:digit2}
\end{align}
so $d_2$ is constant on each interval
$\left[\frac{k}{100}, \frac{k+1}{100}\right)$.

\end{enumerate}


% Unique link targets for equations, figures and tables. Numbers restart in
% every item (\numberwithin{...}{enumi} in main.tex), so without the
% subsection in the name, e.g. equation 1 of item 1 in Formulae and in
% Problems would get the same hyperref target and links could jump to the
% wrong place ("destination with the same identifier" warnings).
\renewcommand{\theHequation}{\thesubsection.\arabic{enumi}.\arabic{equation}}
\renewcommand{\theHfigure}{\thesubsection.\arabic{enumi}.\arabic{figure}}
\renewcommand{\theHtable}{\thesubsection.\arabic{enumi}.\arabic{table}}

\begin{enumerate}[label=\thesubsection.\arabic*, ref=\thesubsection.\theenumi]

%---------------------------------------------------------------------
\item The Z-transform of the PMF of an integer-valued random variable $X$
is defined as
\begin{align}
M_X(z) = \sum_{k=-\infty}^{\infty} p_X(k)\, z^{-k} = \mean{z^{-X}}
\label{eq:ztrans-def}
\end{align}

%---------------------------------------------------------------------
\item If $X_1, \dots, X_n$ are independent and $Y = X_1 + \dots + X_n$, then
\begin{align}
M_Y(z) = \prod_{i=1}^{n} M_{X_i}(z)
\label{eq:ztrans-prod}
\end{align}

%---------------------------------------------------------------------
\item For a Bernoulli random variable $X$ with $\pr{X = 1} = p$ and
$\pr{X = 0} = q = 1 - p$,
\begin{align}
M_X(z) = q + p\, z^{-1}
\label{eq:ztrans-bern}
\end{align}

%---------------------------------------------------------------------
\item If $Y$ is the sum of $n$ independent Bernoulli random variables, each
with parameter $p$, then
\begin{align}
p_Y(k) = \binom{n}{k} p^k q^{n-k}, \quad k = 0, 1, \dots, n
\label{eq:binom-pmf}
\end{align}

%---------------------------------------------------------------------
\item The joint Z-transform of two integer-valued random variables $X$ and
$Y$ is
\begin{align}
M_{X,Y}\brak{z_1, z_2} = \mean{z_1^{-X} z_2^{-Y}}
= \sum_{j}\sum_{k} \pr{X = j, Y = k}\, z_1^{-j} z_2^{-k}
\label{eq:ztrans-joint}
\end{align}
so $\pr{X = j, Y = k}$ is the coefficient of $z_1^{-j} z_2^{-k}$. If
$\brak{X, Y} = \sum_i \brak{X_i, Y_i}$ with the pairs $\brak{X_i, Y_i}$
independent, then
\begin{align}
M_{X,Y}\brak{z_1, z_2} = \prod_{i} M_{X_i,Y_i}\brak{z_1, z_2}
\label{eq:ztrans-joint-prod}
\end{align}

%---------------------------------------------------------------------
\item For a fair die $X$, uniform on $\cbrak{1, \dots, 6}$,
\begin{align}
M_X(z) = \frac{1}{6}\sum_{k=1}^{6} z^{-k}
\label{eq:ztrans-die}
\end{align}
If $S = X_1 + X_2$ for two independent fair dice,
\begin{align}
p_S(k) = \frac{6 - \abs{k - 7}}{36}, \quad k = 2, \dots, 12
\label{eq:dice-sum-pmf}
\end{align}

%---------------------------------------------------------------------
\item A random variable $X$ is Gaussian (normal), $X \sim \gauss{\mu}{\sigma^2}$,
if its PDF is
\begin{align}
f(x) = \frac{1}{\sigma\sqrt{2\pi}}
\exp\brak{-\frac{\brak{x - \mu}^2}{2\sigma^2}}, \quad
x \in \brak{-\infty, \infty}
\label{eq:gauss-pdf}
\end{align}

%---------------------------------------------------------------------
\item Other standard distributions:
\begin{enumerate}[label=\alph*)]
\item An exponential random variable with rate $\lambda > 0$ is continuous,
takes only non-negative values, and has PDF
\begin{align}
f(x) =
\begin{cases}
\lambda e^{-\lambda x} & x \ge 0 \\
0 & x < 0
\end{cases}
\label{eq:exp-pdf}
\end{align}
\item A Poisson random variable $X$ with parameter $\lambda > 0$ is
discrete. It counts the number of events in a fixed interval when events
occur independently at an average rate $\lambda$, so it takes only the
values $0, 1, 2, \dots$, with PMF
\begin{align}
\pr{X = k} = \frac{\lambda^{k} e^{-\lambda}}{k!}, \quad k = 0, 1, 2, \dots
\label{eq:poisson-pmf}
\end{align}
Being discrete, it has a PMF, not a PDF over the real line.
\item A uniform random variable on $[a, b]$ is continuous with PDF
\begin{align}
f(x) =
\begin{cases}
\frac{1}{b - a} & a \le x \le b \\
0 & \text{otherwise}
\end{cases}
\label{eq:unif-pdf}
\end{align}
which is constant on $[a, b]$ and zero outside it.
\end{enumerate}

%---------------------------------------------------------------------
\item Conditional probability:
\begin{align}
\pr{A \mid B} = \frac{\pr{AB}}{\pr{B}}
\label{eq:cond-prob}
\end{align}

%---------------------------------------------------------------------
\item If $g$ is an odd function, $g(-x) = -g(x)$, then
\begin{align}
\int_{-a}^{a} g(x)\, dx = 0
\label{eq:odd-int}
\end{align}

%---------------------------------------------------------------------
\item For $m \neq -1$,
\begin{align}
\int x^{m}\, dx = \frac{x^{m+1}}{m+1} + C
\label{eq:power-rule}
\end{align}

%---------------------------------------------------------------------
\item The sum of the first $n$ natural numbers is
\begin{align}
\sum_{k=1}^{n} k = \frac{n(n+1)}{2}
\label{eq:arith-sum}
\end{align}

%---------------------------------------------------------------------
\item Asymptotic notation. For positive functions $f(n)$ and $g(n)$,
\begin{align}
f(n) = o\brak{g(n)} &\iff \lim_{n \to \infty} \frac{f(n)}{g(n)} = 0
\label{eq:little-o} \\
f(n) = O\brak{g(n)} &\iff f(n) \le c\, g(n) \text{ for large } n
\label{eq:big-o} \\
f(n) = \Theta\brak{g(n)} &\iff c_1 g(n) \le f(n) \le c_2 g(n)
\text{ for large } n
\label{eq:theta}
\end{align}
for some constants $c, c_1, c_2 > 0$. In words, $o$ (little-o) means $f$ grows \emph{strictly slower} than $g$:
however large a multiple of $f$ is taken, $g$ eventually overtakes it.
$O$ (big-O) means $f$ grows \emph{no faster} than $g$, possibly at the
same rate. $\Theta$ means $f$ and $g$ grow at the \emph{same} rate, up to
constant factors. So $o$ is the strict version of $O$, just as $<$ is the
strict version of $\le$.

%---------------------------------------------------------------------
\item For any constants $\epsilon > 0$ and $0 \le a < b$,
\begin{align}
\log(n) = o\brak{n^{\epsilon}}, \qquad n^{a} = o\brak{n^{b}}
\label{eq:growth}
\end{align}

%---------------------------------------------------------------------
\item
\begin{align}
\log(n!) = \Theta\brak{n \log(n)}
\label{eq:log-fact}
\end{align}

%---------------------------------------------------------------------
\item Master theorem. For $T(n) = a\,T\brak{\frac{n}{b}} + f(n)$ with
$a \ge 1$, $b > 1$,
\begin{align}
T(n) =
\begin{cases}
\Theta\brak{n^{\log_b a}} & f(n) = O\brak{n^{\log_b a - \epsilon}}
\text{ for some } \epsilon > 0 \\
\Theta\brak{n^{\log_b a}\log(n)} & f(n) = \Theta\brak{n^{\log_b a}}
\end{cases}
\label{eq:master}
\end{align}

%---------------------------------------------------------------------
\item The (one-sided) Z-transform of a sequence $x[k]$, $k \ge 0$, is
\begin{align}
X(z) = \sum_{k=0}^{\infty} x[k]\, z^{-k}
\label{eq:zseq-def}
\end{align}
If $x[-1] = 0$, a delay by one step multiplies the transform by $z^{-1}$:
\begin{align}
x[k-1] \system{Z} z^{-1} X(z)
\label{eq:zseq-shift}
\end{align}

%---------------------------------------------------------------------
\item For a constant $a$ and $\abs{z} > \abs{a}$,
\begin{align}
a^{k} &\system{Z} \frac{1}{1 - a z^{-1}}
\label{eq:zseq-geo} \\
(k+1)\, a^{k} &\system{Z} \frac{1}{\brak{1 - a z^{-1}}^{2}}
\label{eq:zseq-geo2} \\
\frac{(k+1)(k+2)}{2} &\system{Z} \frac{1}{\brak{1 - z^{-1}}^{3}}
\label{eq:zseq-geo3}
\end{align}

%---------------------------------------------------------------------
\item A relation $R$ on a set $A$ is
\begin{align}
\text{reflexive if } & (x, x) \in R \ \ \forall x \in A
\label{eq:rel-ref} \\
\text{symmetric if } & (x, y) \in R \implies (y, x) \in R
\label{eq:rel-sym} \\
\text{transitive if } & (x, y), (y, z) \in R \implies (x, z) \in R
\label{eq:rel-trans} \\
\text{antisymmetric if } & (x, y), (y, x) \in R \implies x = y
\label{eq:rel-anti}
\end{align}

%---------------------------------------------------------------------
\item The characteristic equation for a matrix $\vec{A}$ is
\begin{align}
f\brak{\lambda} = \mydet{\vec{A} - \lambda\vec{I}} = 0
\label{eq:chareq}
\end{align}

%---------------------------------------------------------------------
\item Cayley-Hamilton theorem:
\begin{align}
f\brak{\lambda} = f\brak{\vec{A}} = \vec{0}
\label{eq:cayley}
\end{align}

%---------------------------------------------------------------------
\item Solving $\vec{A}\vec{x} = \vec{c}$ by row reduction of the augmented
matrix $\myvec{\vec{A} & \vec{c}}$. Elementary row operations, such as
\begin{align}
R_2 \leftarrow \alpha R_2 + \beta R_1, \quad \alpha \ne 0
\label{eq:row-op}
\end{align}
do not change the solutions. After reduction:
\begin{enumerate}[label=\alph*)]
\item a row $\myvec{0 & \dots & 0 & | & 0}$ leaves a variable free, giving
infinitely many solutions;
\item a row $\myvec{0 & \dots & 0 & | & d}$ with $d \ne 0$ means there is
no solution.
\end{enumerate}
\label{fm:row-reduce}

%---------------------------------------------------------------------
\item For an $n \times n$ matrix $\vec{A}$ and a scalar $c$,
\begin{align}
\det\brak{c\vec{A}} = c^{n}\det\brak{\vec{A}}
\label{eq:det-scale}
\end{align}

%---------------------------------------------------------------------
\item In $n$-bit two's complement, the bits $b_{n-1} \dots b_1 b_0$
represent
\begin{align}
x = -b_{n-1} 2^{n-1} + \sum_{i=0}^{n-2} b_i 2^{i}
\label{eq:twos-val}
\end{align}
and the representable range is
\begin{align}
-2^{n-1} \le x \le 2^{n-1} - 1
\label{eq:twos-range}
\end{align}
An arithmetic operation overflows if its true result lies outside
\eqref{eq:twos-range}.

%---------------------------------------------------------------------
\item IEEE 754 single precision. A 32-bit word with sign bit $s$, 8-bit
exponent $E$ and 23-bit fraction $f$ represents (for $0 < E < 255$)
\begin{align}
x = \brak{-1}^{s} \times \brak{1.f}_2 \times 2^{E - 127}
\label{eq:ieee}
\end{align}

%---------------------------------------------------------------------
\item Cyclic Redundancy Check. For data bits $D(x)$ and a generator $G(x)$
of degree $r$, the check bits appended to the data are
\begin{align}
R(x) = x^{r} D(x) \bmod G(x)
\label{eq:crc}
\end{align}
computed with modulo-2 (XOR) arithmetic. In bits: append $r$ zeros to
the data and divide by the generator pattern.

%---------------------------------------------------------------------
\item Boolean algebra:
\begin{align}
X + X' &= 1 \label{eq:bool-comp} \\
XY + X'Y &= Y \label{eq:bool-comb}
\end{align}

%---------------------------------------------------------------------
\item An active-high $2:4$ decoder with inputs $A$ (MSB) and $B$ gives
\begin{align}
D_0 = A'B', \quad D_1 = A'B, \quad D_2 = AB', \quad D_3 = AB
\label{eq:decoder}
\end{align}
A $4$-to-$1$ multiplexer with select lines $S_1 S_0$ and inputs
$I_0, \dots, I_3$ gives
\begin{align}
Y = S_1'S_0' I_0 + S_1'S_0 I_1 + S_1 S_0' I_2 + S_1 S_0 I_3
\label{eq:mux}
\end{align}

%---------------------------------------------------------------------
\item The Jaccard coefficient of two sets $A$, $B$ is
\begin{align}
J(A, B) = \frac{\abs{A \cap B}}{\abs{A \cup B}}
\label{eq:jaccard}
\end{align}
A spanning tree of a connected graph with $n$ vertices has $n - 1$ edges.
The complete graph $K_n$ has $\frac{n(n-1)}{2}$ edges.

%---------------------------------------------------------------------
\item Hashing with chaining. A key $k$ is placed in the chain of slot
\begin{align}
h(k) = k \bmod m
\label{eq:hash}
\end{align}
for a table of $m$ slots. A chain's length is the number of keys that
hash to its slot.

%---------------------------------------------------------------------
\item Merging sorted arrays $L$ and $R$ outputs $L$ followed by $R$ if and
only if
\begin{align}
\max(L) \le \min(R)
\label{eq:merge-void}
\end{align}

%---------------------------------------------------------------------
\item A stack is last-in first-out: elements are popped in the reverse of
the order in which they were pushed. A queue is first-in first-out:
elements are dequeued in the order in which they were enqueued.
\label{fm:stack-queue}

%---------------------------------------------------------------------
\item Binary semaphore $S$:
\begin{align}
\text{Wait}(S) &: \text{block until } S = 1, \text{ then set } S = 0
\label{eq:sem-wait} \\
\text{Signal}(S) &: \text{set } S = 1 \text{ (waking one waiting process)}
\label{eq:sem-signal}
\end{align}
A binary semaphore initialised to $1$, with Wait before and Signal after a
block of code, lets only one process run that block at a time (mutual
exclusion).

%---------------------------------------------------------------------
\item C scope rules. A declaration in an inner block hides a declaration
of the same name in an outer scope. The same name cannot be declared twice
in the same scope. For any variable \texttt{x},
\begin{align}
\texttt{*(\&x)} = \texttt{x}
\label{eq:deref}
\end{align}
\label{fm:c-scope}

%---------------------------------------------------------------------
\item The second digit after the decimal point of $r \in (0, 1)$ is
\begin{align}
d_2(r) = \lfloor 100 r \rfloor \bmod 10
\label{eq:digit2}
\end{align}
so $d_2$ is constant on each interval
$\left[\frac{k}{100}, \frac{k+1}{100}\right)$.

\end{enumerate}


% Unique link targets for equations, figures and tables. Numbers restart in
% every item (\numberwithin{...}{enumi} in main.tex), so without the
% subsection in the name, e.g. equation 1 of item 1 in Formulae and in
% Problems would get the same hyperref target and links could jump to the
% wrong place ("destination with the same identifier" warnings).
\renewcommand{\theHequation}{\thesubsection.\arabic{enumi}.\arabic{equation}}
\renewcommand{\theHfigure}{\thesubsection.\arabic{enumi}.\arabic{figure}}
\renewcommand{\theHtable}{\thesubsection.\arabic{enumi}.\arabic{table}}

\begin{enumerate}[label=\thesubsection.\arabic*, ref=\thesubsection.\theenumi]

%---------------------------------------------------------------------
\item The Z-transform of the PMF of an integer-valued random variable $X$
is defined as
\begin{align}
M_X(z) = \sum_{k=-\infty}^{\infty} p_X(k)\, z^{-k} = \mean{z^{-X}}
\label{eq:ztrans-def}
\end{align}

%---------------------------------------------------------------------
\item If $X_1, \dots, X_n$ are independent and $Y = X_1 + \dots + X_n$, then
\begin{align}
M_Y(z) = \prod_{i=1}^{n} M_{X_i}(z)
\label{eq:ztrans-prod}
\end{align}

%---------------------------------------------------------------------
\item For a Bernoulli random variable $X$ with $\pr{X = 1} = p$ and
$\pr{X = 0} = q = 1 - p$,
\begin{align}
M_X(z) = q + p\, z^{-1}
\label{eq:ztrans-bern}
\end{align}

%---------------------------------------------------------------------
\item If $Y$ is the sum of $n$ independent Bernoulli random variables, each
with parameter $p$, then
\begin{align}
p_Y(k) = \binom{n}{k} p^k q^{n-k}, \quad k = 0, 1, \dots, n
\label{eq:binom-pmf}
\end{align}

%---------------------------------------------------------------------
\item The joint Z-transform of two integer-valued random variables $X$ and
$Y$ is
\begin{align}
M_{X,Y}\brak{z_1, z_2} = \mean{z_1^{-X} z_2^{-Y}}
= \sum_{j}\sum_{k} \pr{X = j, Y = k}\, z_1^{-j} z_2^{-k}
\label{eq:ztrans-joint}
\end{align}
so $\pr{X = j, Y = k}$ is the coefficient of $z_1^{-j} z_2^{-k}$. If
$\brak{X, Y} = \sum_i \brak{X_i, Y_i}$ with the pairs $\brak{X_i, Y_i}$
independent, then
\begin{align}
M_{X,Y}\brak{z_1, z_2} = \prod_{i} M_{X_i,Y_i}\brak{z_1, z_2}
\label{eq:ztrans-joint-prod}
\end{align}

%---------------------------------------------------------------------
\item For a fair die $X$, uniform on $\cbrak{1, \dots, 6}$,
\begin{align}
M_X(z) = \frac{1}{6}\sum_{k=1}^{6} z^{-k}
\label{eq:ztrans-die}
\end{align}
If $S = X_1 + X_2$ for two independent fair dice,
\begin{align}
p_S(k) = \frac{6 - \abs{k - 7}}{36}, \quad k = 2, \dots, 12
\label{eq:dice-sum-pmf}
\end{align}

%---------------------------------------------------------------------
\item A random variable $X$ is Gaussian (normal), $X \sim \gauss{\mu}{\sigma^2}$,
if its PDF is
\begin{align}
f(x) = \frac{1}{\sigma\sqrt{2\pi}}
\exp\brak{-\frac{\brak{x - \mu}^2}{2\sigma^2}}, \quad
x \in \brak{-\infty, \infty}
\label{eq:gauss-pdf}
\end{align}

%---------------------------------------------------------------------
\item Other standard distributions:
\begin{enumerate}[label=\alph*)]
\item An exponential random variable with rate $\lambda > 0$ is continuous,
takes only non-negative values, and has PDF
\begin{align}
f(x) =
\begin{cases}
\lambda e^{-\lambda x} & x \ge 0 \\
0 & x < 0
\end{cases}
\label{eq:exp-pdf}
\end{align}
\item A Poisson random variable $X$ with parameter $\lambda > 0$ is
discrete. It counts the number of events in a fixed interval when events
occur independently at an average rate $\lambda$, so it takes only the
values $0, 1, 2, \dots$, with PMF
\begin{align}
\pr{X = k} = \frac{\lambda^{k} e^{-\lambda}}{k!}, \quad k = 0, 1, 2, \dots
\label{eq:poisson-pmf}
\end{align}
Being discrete, it has a PMF, not a PDF over the real line.
\item A uniform random variable on $[a, b]$ is continuous with PDF
\begin{align}
f(x) =
\begin{cases}
\frac{1}{b - a} & a \le x \le b \\
0 & \text{otherwise}
\end{cases}
\label{eq:unif-pdf}
\end{align}
which is constant on $[a, b]$ and zero outside it.
\end{enumerate}

%---------------------------------------------------------------------
\item Conditional probability:
\begin{align}
\pr{A \mid B} = \frac{\pr{AB}}{\pr{B}}
\label{eq:cond-prob}
\end{align}

%---------------------------------------------------------------------
\item If $g$ is an odd function, $g(-x) = -g(x)$, then
\begin{align}
\int_{-a}^{a} g(x)\, dx = 0
\label{eq:odd-int}
\end{align}

%---------------------------------------------------------------------
\item For $m \neq -1$,
\begin{align}
\int x^{m}\, dx = \frac{x^{m+1}}{m+1} + C
\label{eq:power-rule}
\end{align}

%---------------------------------------------------------------------
\item The sum of the first $n$ natural numbers is
\begin{align}
\sum_{k=1}^{n} k = \frac{n(n+1)}{2}
\label{eq:arith-sum}
\end{align}

%---------------------------------------------------------------------
\item Asymptotic notation. For positive functions $f(n)$ and $g(n)$,
\begin{align}
f(n) = o\brak{g(n)} &\iff \lim_{n \to \infty} \frac{f(n)}{g(n)} = 0
\label{eq:little-o} \\
f(n) = O\brak{g(n)} &\iff f(n) \le c\, g(n) \text{ for large } n
\label{eq:big-o} \\
f(n) = \Theta\brak{g(n)} &\iff c_1 g(n) \le f(n) \le c_2 g(n)
\text{ for large } n
\label{eq:theta}
\end{align}
for some constants $c, c_1, c_2 > 0$. In words, $o$ (little-o) means $f$ grows \emph{strictly slower} than $g$:
however large a multiple of $f$ is taken, $g$ eventually overtakes it.
$O$ (big-O) means $f$ grows \emph{no faster} than $g$, possibly at the
same rate. $\Theta$ means $f$ and $g$ grow at the \emph{same} rate, up to
constant factors. So $o$ is the strict version of $O$, just as $<$ is the
strict version of $\le$.

%---------------------------------------------------------------------
\item For any constants $\epsilon > 0$ and $0 \le a < b$,
\begin{align}
\log(n) = o\brak{n^{\epsilon}}, \qquad n^{a} = o\brak{n^{b}}
\label{eq:growth}
\end{align}

%---------------------------------------------------------------------
\item
\begin{align}
\log(n!) = \Theta\brak{n \log(n)}
\label{eq:log-fact}
\end{align}

%---------------------------------------------------------------------
\item Master theorem. For $T(n) = a\,T\brak{\frac{n}{b}} + f(n)$ with
$a \ge 1$, $b > 1$,
\begin{align}
T(n) =
\begin{cases}
\Theta\brak{n^{\log_b a}} & f(n) = O\brak{n^{\log_b a - \epsilon}}
\text{ for some } \epsilon > 0 \\
\Theta\brak{n^{\log_b a}\log(n)} & f(n) = \Theta\brak{n^{\log_b a}}
\end{cases}
\label{eq:master}
\end{align}

%---------------------------------------------------------------------
\item The (one-sided) Z-transform of a sequence $x[k]$, $k \ge 0$, is
\begin{align}
X(z) = \sum_{k=0}^{\infty} x[k]\, z^{-k}
\label{eq:zseq-def}
\end{align}
If $x[-1] = 0$, a delay by one step multiplies the transform by $z^{-1}$:
\begin{align}
x[k-1] \system{Z} z^{-1} X(z)
\label{eq:zseq-shift}
\end{align}

%---------------------------------------------------------------------
\item For a constant $a$ and $\abs{z} > \abs{a}$,
\begin{align}
a^{k} &\system{Z} \frac{1}{1 - a z^{-1}}
\label{eq:zseq-geo} \\
(k+1)\, a^{k} &\system{Z} \frac{1}{\brak{1 - a z^{-1}}^{2}}
\label{eq:zseq-geo2} \\
\frac{(k+1)(k+2)}{2} &\system{Z} \frac{1}{\brak{1 - z^{-1}}^{3}}
\label{eq:zseq-geo3}
\end{align}

%---------------------------------------------------------------------
\item A relation $R$ on a set $A$ is
\begin{align}
\text{reflexive if } & (x, x) \in R \ \ \forall x \in A
\label{eq:rel-ref} \\
\text{symmetric if } & (x, y) \in R \implies (y, x) \in R
\label{eq:rel-sym} \\
\text{transitive if } & (x, y), (y, z) \in R \implies (x, z) \in R
\label{eq:rel-trans} \\
\text{antisymmetric if } & (x, y), (y, x) \in R \implies x = y
\label{eq:rel-anti}
\end{align}

%---------------------------------------------------------------------
\item The characteristic equation for a matrix $\vec{A}$ is
\begin{align}
f\brak{\lambda} = \mydet{\vec{A} - \lambda\vec{I}} = 0
\label{eq:chareq}
\end{align}

%---------------------------------------------------------------------
\item Cayley-Hamilton theorem:
\begin{align}
f\brak{\lambda} = f\brak{\vec{A}} = \vec{0}
\label{eq:cayley}
\end{align}

%---------------------------------------------------------------------
\item Solving $\vec{A}\vec{x} = \vec{c}$ by row reduction of the augmented
matrix $\myvec{\vec{A} & \vec{c}}$. Elementary row operations, such as
\begin{align}
R_2 \leftarrow \alpha R_2 + \beta R_1, \quad \alpha \ne 0
\label{eq:row-op}
\end{align}
do not change the solutions. After reduction:
\begin{enumerate}[label=\alph*)]
\item a row $\myvec{0 & \dots & 0 & | & 0}$ leaves a variable free, giving
infinitely many solutions;
\item a row $\myvec{0 & \dots & 0 & | & d}$ with $d \ne 0$ means there is
no solution.
\end{enumerate}
\label{fm:row-reduce}

%---------------------------------------------------------------------
\item For an $n \times n$ matrix $\vec{A}$ and a scalar $c$,
\begin{align}
\det\brak{c\vec{A}} = c^{n}\det\brak{\vec{A}}
\label{eq:det-scale}
\end{align}

%---------------------------------------------------------------------
\item In $n$-bit two's complement, the bits $b_{n-1} \dots b_1 b_0$
represent
\begin{align}
x = -b_{n-1} 2^{n-1} + \sum_{i=0}^{n-2} b_i 2^{i}
\label{eq:twos-val}
\end{align}
and the representable range is
\begin{align}
-2^{n-1} \le x \le 2^{n-1} - 1
\label{eq:twos-range}
\end{align}
An arithmetic operation overflows if its true result lies outside
\eqref{eq:twos-range}.

%---------------------------------------------------------------------
\item IEEE 754 single precision. A 32-bit word with sign bit $s$, 8-bit
exponent $E$ and 23-bit fraction $f$ represents (for $0 < E < 255$)
\begin{align}
x = \brak{-1}^{s} \times \brak{1.f}_2 \times 2^{E - 127}
\label{eq:ieee}
\end{align}

%---------------------------------------------------------------------
\item Cyclic Redundancy Check. For data bits $D(x)$ and a generator $G(x)$
of degree $r$, the check bits appended to the data are
\begin{align}
R(x) = x^{r} D(x) \bmod G(x)
\label{eq:crc}
\end{align}
computed with modulo-2 (XOR) arithmetic. In bits: append $r$ zeros to
the data and divide by the generator pattern.

%---------------------------------------------------------------------
\item Boolean algebra:
\begin{align}
X + X' &= 1 \label{eq:bool-comp} \\
XY + X'Y &= Y \label{eq:bool-comb}
\end{align}

%---------------------------------------------------------------------
\item An active-high $2:4$ decoder with inputs $A$ (MSB) and $B$ gives
\begin{align}
D_0 = A'B', \quad D_1 = A'B, \quad D_2 = AB', \quad D_3 = AB
\label{eq:decoder}
\end{align}
A $4$-to-$1$ multiplexer with select lines $S_1 S_0$ and inputs
$I_0, \dots, I_3$ gives
\begin{align}
Y = S_1'S_0' I_0 + S_1'S_0 I_1 + S_1 S_0' I_2 + S_1 S_0 I_3
\label{eq:mux}
\end{align}

%---------------------------------------------------------------------
\item The Jaccard coefficient of two sets $A$, $B$ is
\begin{align}
J(A, B) = \frac{\abs{A \cap B}}{\abs{A \cup B}}
\label{eq:jaccard}
\end{align}
A spanning tree of a connected graph with $n$ vertices has $n - 1$ edges.
The complete graph $K_n$ has $\frac{n(n-1)}{2}$ edges.

%---------------------------------------------------------------------
\item Hashing with chaining. A key $k$ is placed in the chain of slot
\begin{align}
h(k) = k \bmod m
\label{eq:hash}
\end{align}
for a table of $m$ slots. A chain's length is the number of keys that
hash to its slot.

%---------------------------------------------------------------------
\item Merging sorted arrays $L$ and $R$ outputs $L$ followed by $R$ if and
only if
\begin{align}
\max(L) \le \min(R)
\label{eq:merge-void}
\end{align}

%---------------------------------------------------------------------
\item A stack is last-in first-out: elements are popped in the reverse of
the order in which they were pushed. A queue is first-in first-out:
elements are dequeued in the order in which they were enqueued.
\label{fm:stack-queue}

%---------------------------------------------------------------------
\item Binary semaphore $S$:
\begin{align}
\text{Wait}(S) &: \text{block until } S = 1, \text{ then set } S = 0
\label{eq:sem-wait} \\
\text{Signal}(S) &: \text{set } S = 1 \text{ (waking one waiting process)}
\label{eq:sem-signal}
\end{align}
A binary semaphore initialised to $1$, with Wait before and Signal after a
block of code, lets only one process run that block at a time (mutual
exclusion).

%---------------------------------------------------------------------
\item C scope rules. A declaration in an inner block hides a declaration
of the same name in an outer scope. The same name cannot be declared twice
in the same scope. For any variable \texttt{x},
\begin{align}
\texttt{*(\&x)} = \texttt{x}
\label{eq:deref}
\end{align}
\label{fm:c-scope}

%---------------------------------------------------------------------
\item The second digit after the decimal point of $r \in (0, 1)$ is
\begin{align}
d_2(r) = \lfloor 100 r \rfloor \bmod 10
\label{eq:digit2}
\end{align}
so $d_2$ is constant on each interval
$\left[\frac{k}{100}, \frac{k+1}{100}\right)$.

\end{enumerate}
